\documentclass[10pt]{article}

\usepackage[empty]{fullpage}
\usepackage{hyperref}
\usepackage{tikz}
\usepackage{soul}

\begin{document}
\noindent
%
% REQUIRED SYLLABUS CONTENT AS OF SPRING 2025:
% 1. Course Information
%     - Course prefix and number, section number, title, and credit hours
%     - Course description (include full Academic Catalogue description with prerequisites)
%     - Semester and year
%     - Class meeting days, time, and location
%     - Faculty member’s name, contact information, and office hours
%     - Required course materials and texts
%     - Assignment descriptions and calculation of grades
%     - Course schedule
%     - Policies governing late work, make-up assignments, and attendance
% 2. Student Learning Information
%     - Course’s relationship to degree program
%     - Student learning outcomes for major/standalone minor
%     - Course learning outcomes (see template)
% 3. All syllabi must contain the following institutional policies. The wording for these policies should be copied and pasted from the syllabus template.
%     - Grading Scale
%     - Honor Code
%     - Policies for use of AI on assignments
%     - Accommodations


\begin{center}
\includegraphics[scale=0.6]{../../../HSC.png} 
\bigskip

\textbf{COMS 261: Computer Science I (3 credits)} \\
Fall 2026
\end{center}

\noindent
\begin{tabular}{|l|l|}
\hline
Instructor & Brian Lins \\ \hline
Email Address & blins@hsc.edu \\ \hline
Course Meeting Times & MWF 10:30 - 11:20am \\ \hline
Course Meeting Location & Pauley 100 \\ \hline
Office Hours & MW 2:30 - 3:30pm \& R 1:00 - 3:00pm (and by appointment) \\ 
& See the course website: \url{https://bclins.github.io} ~ \\ \hline
Textbook & Think Python, 2nd edition by Allen Downey \\ \hline
\end{tabular}

\subsubsection*{Course Description}

Discussion of algorithms, programs, and computers. Extensive work in the preparation, running, debugging, and documenting of programs. Problem-solving is emphasized.

\subsubsection*{Course Learning Objectives}


By the end of this course, students will be able to:

\begin{itemize}
\item Read and write programs in the Python programming language.
\item Apply algorithmic thinking to solve problems.  
\item Explain fundamental concepts of computing including abstraction, algorithms, and data representation. 
%%%% Williams CLO's:
%\item Write, run, and debug Python programs using variables, conditionals, loops, functions, and data structures.
%\item Apply algorithmic thinking to solve problems in contexts such as encryption, games, and data manipulation.
%\item Use the Linux command line for file management, scripting, and interacting with development tools.
%\item Explain and apply fundamental concepts of computing, including abstraction, algorithms, and data representation.
%\item Document and reflect on programming projects through a personal portfolio website.
%\item Evaluate how computing intersects with issues of privacy, security, and society.
\end{itemize}

\subsubsection*{Required Materials}

None. See the course website for links to the free textbook.

%\subsubsection*{Assignment Descriptions and Grade}
%
\subsubsection*{Attendance Policy}

Attendance in this class is required. Repeated absences may result in a forced withdrawal from the course. You are responsible for any material you miss due to absence. Please let me know ahead of time if you know that you will not be able to attend class.

\subsubsection*{Grading Policy}

The term grade will be based on the following factors.

\begin{center}
\begin{tabular}{|l|c|}
\hline
Component      & Proportion \\ \hline
Midterms  & 60\% \\
Projects  & 10\% \\
Final Exam  & 30\% \\ \hline
\end{tabular}
\end{center}


\subsection*{In-Class Labs}
 
Programming is something that you can only learn by doing. During most class periods you will be asked to write short pieces of code. If you aren't able to finish a task, then you should try to figure it out later on your own or come to office hours for help. These in-class programming exercises won't be collected, although similar questions are likely to appear on tests and projects.


\subsection*{Projects}

There will be several programming projects throughout the semester.  You must complete these projects on your own.  Do not ask AI or other students in this class for help on projects. If you include any code from another student or from the internet, that is plagiarism and will be treated as such. To avoid plagiarizing other students, you should avoid looking at other students' code because you cannot ``unsee" it.  You should also never share your code with other students. For the same reason, you should not ask AI for help, not even with small parts of a programming assignment.  If you need help with a project, see me or one of the student tutors for this course.  

Projects will be graded based on several factors.  Code should be written following the Python PEP-8 (\verb|https://peps.python.org/pep-0008/|) style guide with clear comments.  It should make sense, and it should run correctly on all test cases.  These factors will be graded holistically using the following rubric.

\begin{itemize}
\item \textbf{Grade: A.} Every part of the assignment is complete and your code is clear and runs correctly.  
\item \textbf{Grade: B.} Almost every part of the assignment is complete and your work is clear. There may be a few minor mistakes, but no major errors.  
\item \textbf{Grade: C.} Most of the assignment is complete, but is either not clear or has significant mistakes.
\item \textbf{Grade: D.} You submit something, but it has major errors or omissions.  
\end{itemize}
 
\noindent
All projects will be posted at least 4 days before they are due.  


\subsection*{Exams}

There will be three in-class midterm exams and a cumulative final. These exams will be announced in advance, and you will know exactly what concepts will be covered on each exam.  One of the midterm exams may be a practical exam where you meet with me one-on-one and demonstrate your coding abilities. 

\subsubsection*{Course Schedule} 

The schedule below is tentative, and may be subject to change. Changes will be announced in class, and you are responsible for knowing about any changes even if you miss the class when they are announced. 

\begin{center}
\begin{tabular}{|c|l|l|}
\hline
Week  & Topic & Projects \\ \hline
1  & Variables, expressions, \& statements & \\
2  & Functions                             & Project 1 \\
3  & Conditionals                          & \\
4  & Recursion                             & \\
5  & More on functions                     & Project 2 \\
6  & Iteration, Midterm 1                  & \\
7  & Strings                               & \\
8  & Lists                                 & Project 3 \\
9  & Dictionaries                          & \\
10 & Tuples                                & \\
11 & File                                  & Project 4 \\
12 & Classes and objects                   & \\
13 & Classes and functions, Midterm 2      & \\
14 & Classes and methods                   & Project 5 \\
15 & Practical exams                       & \\ \hline
\end{tabular}
\end{center}

\subsubsection*{Late Work and Make-Up Assignment Policy}

Late projects will be accepted, but with a 10\% grade penalty.  The penalty increases to 20\% for projects that are more than one week late. 

\subsubsection*{Grading Scale} 

This course adheres to the grades and quality points described in the \href{https://www.hsc.edu/academic-catalogues}{Academic Catalogue}. Consult the Academic Catalogue for a detailed description. 


\subsubsection*{Honor Code}

Students are expected to abide by the Honor Code for all assignments unless a professor indicates otherwise. Students should consult the Academic Catalogue and The Key: The Hampden-Sydney College Student Handbook for the College’s description of the Honor Code and what it identifies as infractions of the Honor Code.

\subsubsection*{Artificial Intelligence Policy}

Artificial intelligence can easily solve all of the programming exercises you will encounter in this course. Watching AI solve problems for you will not help you learn computer science. For that reason, AI use is strictly prohibited on projects and tests in this class.  
You are, however, welcome to use AI to help study before exams and to explain material that you didn't understand in class.  
%Artificial intelligence (AI) generators and large language models (LLMs) often rely on existing published materials, and copying or paraphrasing materials generated by AI without attribution is plagiarism. Professors may permit students to use AI generators or LLMs in a variety of ways in their own classes. Those students, however, must not assume that those policies transfer to other classes.

\subsubsection*{Accommodations}

Hampden-Sydney College is committed to ensuring equitable access to its education programs for all students. Under the administration of the Department of Culture and Community, the Office of Accessibility Services (OAS) coordinates reasonable accommodations for qualified students with disabilities. If you wish to seek accommodations for this class, please contact Dr. Melissa Wood, Director of Title IX, Access, and Student Advocacy, at 434-223-6061 or at \url{mwood@hsc.edu}. Additional information may be found here: \url{https://www.hsc.edu/academics/academic-services/disability-services}. Appropriate documentation of disability will be required. For students who have an accommodations letter from OAS, it is essential that you correspond with your professor as soon as possible to discuss your accommodation needs for the course so that appropriate arrangements may be made. 
\end{document}


